% Generated draft: requires mathematical and source review.
\documentclass[11pt]{article}
\usepackage[margin=25mm]{geometry}
\usepackage{fontspec}
\setmainfont{Latin Modern Roman}
\usepackage{amsmath,amssymb,mathtools,unicode-math}
\setmathfont{Latin Modern Math}
\usepackage{enumitem,needspace,xurl,hyperref}
\setlength{\parindent}{0pt}
\setlength{\parskip}{5pt}
\setlength{\emergencystretch}{3em}
\newcommand{\N}{\mathbb{N}}
\newcommand{\Z}{\mathbb{Z}}
\newcommand{\Q}{\mathbb{Q}}
\newcommand{\R}{\mathbb{R}}
\newcommand{\C}{\mathbb{C}}
\newcommand{\F}{\mathbb{F}}
\newcommand{\A}{\mathbb{A}}
\newcommand{\PP}{\mathbb{P}}
\newcommand{\E}{\mathbb{E}}
\newcommand{\eps}{\varepsilon}
\newcommand{\dd}{\,\mathrm{d}}
\begin{document}
\section*{UnsolvedMath ID: 30001046}
\textbf{Source problem code:} catalogue entry
\subsection{Definitions and mathematical statement}
Let \(M\) be a connected finite-volume locally symmetric space of noncompact type, with the rational structure needed to define its \(\Q\)-rank. In an arithmetic realization one writes \(M=\Gamma\backslash\mathbf G(\R)/K\), and \(\operatorname{rank}_{\Q}\mathbf G\) is the dimension of a maximal \(\Q\)-split torus. This is not the real rank of the symmetric space. The extraction does not specify the arithmetic realization, irreducibility, torsion-freeness, or orientability conventions.

For an oriented noncompact \(n\)-manifold, the relevant simplicial volume is the locally finite version
\[
 \|M\|_{\mathrm{lf}}=
 \inf\left\{\sum_i|a_i|:
 c=\sum_i a_i\sigma_i\text{ is a locally finite real singular cycle},
 \ [c]=[M]_{\mathrm{lf}}\right\}\in[0,\infty].
\]
Locally finite means that every compact subset meets the images of only finitely many simplices in the chain. The fundamental class lies in \(H_n^{\mathrm{lf}}(M;\R)\). No finiteness of the infimum is assumed. Nonorientable or orbifold versions require the intended additional convention.

The source asks the two universal implications
\[
 \operatorname{rank}_{\Q}(M)=1\quad\Longrightarrow\quad
 \|M\|_{\mathrm{lf}}\ne0\ ?
\]
\[
 \operatorname{rank}_{\Q}(M)=2\quad\Longrightarrow\quad
 \|M\|_{\mathrm{lf}}=0\ ?
\]
The second question's requested conclusion is vanishing, not positivity; the two directions must not be combined into an unprovided general rank formula. Nonzero in the first question allows infinity unless finiteness is separately established. Ordinary finite-chain top homology on an open manifold, Lipschitz simplicial volume, and relative simplicial volume of a compact core are distinct constructions and are not automatically substituted for \(\|M\|_{\mathrm{lf}}\).

The supplied assessment dated 2026-08-21 reports positivity in large \(\Q\)-rank-one classes but says the general rank-one and rank-two questions are not fully settled. The cited paper labelled real rank two does not by that label alone address every \(\Q\)-rank-two case. Exact excluded symmetric factors, arithmetic and irreducibility hypotheses, and the comparison of simplicial-volume notions must be checked before applying any reported partial theorem.
\subsection{Short English statement}
For the locally finite fundamental-class norm of a finite-volume locally symmetric space, the questions ask for nonvanishing in rational rank one and vanishing in rational rank two. Rational rank and real rank, and ordinary versus Lipschitz or relative simplicial volume, must remain distinct.
\subsection{Sources}
\begin{itemize}
\item Local UnsolvedMath registry, permanent ID \(30001046\), exact source code \texttt{catalogue entry}, proposed year 2008; no individual proposer is supplied.
\item \emph{Topologie (2009)}, Oberwolfach Reports, DOI 10.4171/owr/2008/43, pp. 2419--2476. \url{https://doi.org/10.4171/owr/2008/43}.
\item Kim--Kim, ``Q-rank one,'' \url{https://arxiv.org/abs/1104.4628}.
\item Kim--Kim, ``real rank two,'' \url{https://arxiv.org/abs/1503.02381}.
\item The supplied \texttt{corrected\_verified} note reports checking mathematical content against the official report. The partially solved label and assessment checked 2026-08-21 are supplied dated curation, not independently verified; exact partial-result hypotheses and current status require review.
\end{itemize}
\end{document}
